\section{Introduction}

\subsection{Study Rationale}
Pancreatic ductal adenocarcinoma is the third leading cause of cancer death in
the United States with five-year overall survival below 13 percent, and the
Whipple procedure is its most technically demanding curative-intent operation.
This protocol pairs a precise, auditable, on-premises LLM-directed robotic
Whipple with the RAS(ON) multi-selective inhibitor daraxonrasib.

\draftinstr{Develop the rationale around the three counterfactual scenarios in
which withholding the LLM + robot + medicine combination shortens PFS and OS;
render \texttt{mermaid/fig-19-counterfactual-scenarios.md} as a TikZ figure.
Politely raise the framing questions: how did we ever rely on humans alone for
this treatment; why did we fully trust a prior trial system with so many
opportunities for compounding human error; why would human peer review continue
for increasingly complex and voluminous AI outputs. Source the clinical cascade
from \texttt{inputs/2030-pdac-1min-final-paper/sections/\{introduction,discussion\}.tex}.}

\subsection{Background}
\draftinstr{Summarize: (a) PDAC epidemiology and the Whipple fistula-sepsis
cascade (2030-pdac introduction); (b) daraxonrasib mechanism, FDA Breakthrough
Therapy designation, RASolute 302 / RASolve 301 (cite the five DARAXONRASIB
entries, brief explanation of each); (c) the current state of AI in trials -
prior AI was prediction-only or narrow, 0 generative LLM devices approved, 1,016
to 1,450+ narrow AI/ML authorizations - from
\texttt{research/\{ChatGPT-5.5-Thinking-Extended-18Jun26,Gemini-3.1-Pro-18Jun26\}.md};
(d) the author's established LLM trust, Aug 2024-Jun 2026, from
\texttt{inputs/author\_works.bib}, rendered via
\texttt{mermaid/fig-23-author-trust-timeline.md}, with brief descriptions only.}

\draftinstr{Add a Background subsection ``Physical AI concerns addressed by this
protocol'' covering all eight items (limitations; patient safety; loss of human
workers; single-source software; proprietary-model dependency; non-domestic
open-source LLMs; overly complex workflows; black boxes), each paired with its
mitigation; render \texttt{mermaid/fig-20-physical-ai-concerns.md} as a TikZ
figure and as a full-width table.}

\subsection{Risk/Benefit Assessment}
\subsubsection{Known Potential Risks}
\draftinstr{Enumerate drug risks (daraxonrasib toxicity), device/robotic risks
(malfunction, AI advisory error, cyber-physical failure, vascular injury), and
operative risks (fistula, hemorrhage, conversion). Source robotic risk
categories from \texttt{inputs/21cfr312\_adapt/01\_preamble\_scope\_definitions.tex}.}
\subsubsection{Known Potential Benefits}
\draftinstr{Precise in-window R0 resection, layered hard safety limits with full
audit trail, optimized drug-restart timing; for a low-survival population.}
\subsubsection{Assessment of Potential Risks and Benefits}
\draftinstr{Render \texttt{mermaid/fig-24-risk-benefit-terminal.md} as a TikZ
figure; conclude the benefit-risk balance is favorable for advanced patients,
consistent with the expanded-access ethic (21 CFR 312.84;
\texttt{research/ChatGPT-5.5-Thinking-Extended-19Jun26.md}).}
